\section*{Capítulo \ref{ch:g4:world-physical-map} --- El mapa físico del mundo}

\begin{solution}{exo:g4:world-physical-map:1}
La \omterm{def:g4:world-physical-map:altitude}{altitud} de un lugar es su altura sobre el \omterm{def:g4:world-physical-map:altitude}{nivel del mar}, en metros. Se
mide desde el \omterm{def:g4:world-physical-map:altitude}{nivel del mar}, el nivel medio de la superficie del \omterm{def:g2:water-on-earth:ocean}{mar}.
\end{solution}

\begin{solution}{exo:g4:world-physical-map:2}
(a) Verde oscuro: entre 0 y 200 metros. (b) Amarillo pálido: entre 500 y
1\,000 metros. (c) Casi blanco: más de 4\,000 metros.
\end{solution}

\begin{solution}{exo:g4:world-physical-map:3}
(a) América del Sur. (b) Asia. (c) África. (d) América del Norte.
\end{solution}

\begin{solution}{exo:g4:world-physical-map:4}
La \omterm{def:g1:landscapes-around-us:landform}{montaña} más alta es el Everest (8\,849 metros). El \omterm{def:g1:landscapes-around-us:water}{lago} más grande es el
mar Caspio; el \omterm{def:g1:landscapes-around-us:water}{lago} más profundo es el \omterm{def:g1:landscapes-around-us:water}{lago} Baikal.
\end{solution}

\begin{solution}{exo:g4:world-physical-map:5}
Cuanto más oscuro es el azul, más profundo es el \omterm{def:g2:water-on-earth:ocean}{mar}: el \omterm{def:g2:water-on-earth:ocean}{mar} es poco
profundo cerca de las \omterm{def:g1:landscapes-around-us:water}{costas} (menos de 200 metros) y mucho más profundo en
el centro de los \omterm{def:g2:water-on-earth:ocean}{océanos} (más de 4\,000 metros en muchos lugares).
\end{solution}

\begin{solution}{exo:g4:world-physical-map:6}
El \omterm{def:g4:world-physical-map:physical-map}{mapa físico} muestra la \omterm{def:g4:world-physical-map:altitude}{altitud} de la superficie. En la Antártida y en
Groenlandia la superficie es la parte de arriba de una gruesa capa de
hielo, que está muy por encima del \omterm{def:g4:world-physical-map:altitude}{nivel del mar}, así que tienen los
colores de las tierras altas.
\end{solution}

\begin{solution}{exo:g4:world-physical-map:7}
El fondo del abismo Challenger (casi 11\,000 metros por debajo del \omterm{def:g4:world-physical-map:altitude}{nivel del
mar}); el fondo medio del \omterm{def:g2:water-on-earth:ocean}{océano} (unos 3\,700 metros por debajo); una playa
(unos 0 metros); un pueblo a 1\,500 metros; la cima del Everest
(8\,849 metros).
\end{solution}

\begin{solution}{exo:g4:world-physical-map:8}
Un \omterm{def:g1:landscapes-around-us:water}{río} está formado por muchos arroyos que se juntan, y la gente no se pone
de acuerdo sobre cuál de ellos es su verdadero comienzo. Medido desde una
fuente más lejana, un \omterm{def:g1:landscapes-around-us:water}{río} es más largo. El Nilo y el Amazonas tienen
longitudes tan parecidas (cada uno, entre unos 6\,700 y 6\,800 kilómetros)
que la respuesta depende de esa elección.
\end{solution}

\begin{solution}{exo:g4:world-physical-map:9}
La cuarta parte de 30 millones son siete millones y medio de kilómetros
cuadrados (30 dividido entre 4); la tercera parte, 10 millones (30 dividido
entre 3). El Sáhara, con ocho millones y medio, es más que la cuarta parte
de África, pero menos que la tercera parte.
\end{solution}

\begin{solution}{exo:g4:world-physical-map:10}
$2500 - 1000 = 1500$: el refugio está 1\,500 metros más alto. El refugio,
entre 2\,000 y 4\,000 metros, es marrón; el \omterm{def:g1:landscapes-around-us:water}{lago}, a 1\,000 metros, está en
el límite entre el amarillo pálido (500 a 1\,000) y el marrón claro
(1\,000 a 2\,000): vale cualquiera de los dos.
\end{solution}

\begin{solution}{exo:g4:world-physical-map:11}
10 kilómetros son 10\,000 metros. $10000 - 8849 = 1151$: el avión está más
alto que el Everest, por 1\,151 metros. Sobre el abismo Challenger, está a
10\,000 metros por encima del \omterm{def:g2:water-on-earth:ocean}{mar}, y el fondo de la fosa está a
11\,000 metros por debajo del \omterm{def:g2:water-on-earth:ocean}{mar}: $10000 + 11000 = 21000$ metros por
encima del fondo de la fosa.
\end{solution}
